\documentclass{scrartcl}

\usepackage[utf8]{inputenc}
\usepackage[english]{babel}
\usepackage{amsmath}
\usepackage{amsfonts}
\usepackage{amssymb}
\usepackage{amsthm}
\usepackage{tikz}
\usepackage{hyphenat}
\usepackage{graphicx}
\usepackage{nicefrac}
\usepackage{booktabs}
\usepackage{todonotes}
\usepackage{bm}
\usepackage{BOONDOX-calo} % for nicer calligraphic symbols
\usepackage{authblk}
\usepackage{subcaption}

\usepackage[square]{natbib}
\bibliographystyle{chicago}

\usepackage{geometry}
\usepackage[
    unicode = true, 
    bookmarks = false, 
    breaklinks = true, 
    pdfborder = {0 0 1},
    colorlinks = true
    ]{hyperref}

\geometry{vmargin = 3cm, hmargin = 3cm}

\usetikzlibrary{
    automata, 
    positioning, 
    calc, 
    shapes, 
    arrows, 
    fit
    }

% \tikzset{
%     every loop/.style = {
%         min distance = 1mm, 
%         out = 25, 
%         in = -25, 
%         looseness = 5
%         }
%     }

\definecolor{PastelGreen}{HTML}{6ECB63}
\definecolor{EarthyGreen}{HTML}{588157}
\definecolor{Azure}{RGB}{0, 128, 255}
\definecolor{Fiery}{HTML}{c1121f}

\theoremstyle{definition}
\newtheorem{theorem}{Theorem}
\newtheorem{lemma}{Lemma}
\newtheorem{corollary}{Corollary}
\newtheorem{example}{Example}
\newtheorem{conjecture}{Conjecture}
\newtheorem{proposition}{Proposition}
\newtheorem{definition}{Definition}
\newtheorem{remark}{Remark}
\newtheorem{claim}{Claim}
\newtheorem{observation}{Observation}

\DeclareMathOperator*{\argmax}{arg\,max}
\DeclareMathOperator*{\argmin}{arg\,min}

\renewcommand{\descriptionlabel}[1]{$\triangleright$ \textit{#1}}

\renewcommand{\phi}{\varphi}
\renewcommand{\epsilon}{\varepsilon}
\renewcommand{\theta}{\vartheta}
\newcommand{\len}{\mathsf{len}}
\newcommand{\net}{\mathsf{net}}
\renewcommand{\min}{\mathit{min}}
\renewcommand{\max}{\mathit{max}}
\newcommand{\avg}{\mathit{avg}}
\newcommand{\ev}{\mathbb{E}}
\newcommand{\merge}{M}

\hypersetup{
    linkcolor = blue,
    citecolor = Azure,
    filecolor = blue,
    urlcolor = blue
    }

\tikzset{
	a/.style={
        draw,
        rounded corners=2pt, 
        minimum width=1cm, 
        minimum height=1cm, 
        thick,
        color=black,
        fill=EarthyGreen, 
        text=white,
        font=\Large
        },
    b/.style={
        draw,
        rounded corners=2pt, 
        minimum width=1cm, 
        minimum height=1cm, 
        thick,
        color=black,
        fill=Fiery, 
        text=white,
        font=\Large
        },
    a-agent/.style={
        draw,
        circle, 
        minimum width=2.5em, 
        thick,
        color=black,
        fill=EarthyGreen, 
        text=white,
        font=\Large
        },
    b-agent/.style={
        draw,
        circle, 
        minimum width=2.5em, 
        thick,
        color=black,
        fill=Fiery, 
        text=white,
        font=\Large
        },
    neutral-agent/.style={
        draw,
        circle, 
        minimum width=2.5em, 
        thick,
        color=black,
        fill=gray, 
        text=white,
        font=\Large
        }
    }


\newcommand{\adrian}[1]{
    {
        \hfill
        \begin{minipage}{0.9\textwidth}
            \begin{flushright}
                \color{PastelGreen}
                \footnotesize
                \small 
                {\footnotesize\textsf{\textbf{Adrian}}}\\ 
                \textsf{#1}
            \end{flushright}
        \end{minipage}
    }
    }


\title{
    Social Learning with a Chameleon AI
    }
\author{Adrian \dots}
\date{}




\begin{document}
\maketitle
% \setcounter{tocdepth}{1}
% \tableofcontents

\begin{abstract}
    \noindent
    We study social learning with a sycophantic AI.
\end{abstract}

\section{Introduction}\label{sec:introduction}
We study a model of social learning in which a population of humans 
tries to estimate an unknown quantity by repeatedly averaging opinions 
over a social network, along the lines of the DeGroot model 
\citep{DeGroot1974ReachingConsensus,Golub2010NaiveLearning},
while incorporating the influence of an AI agent. 
The AI agent learns from the population and 
feeds information back   to the humans, though not necessarily in an honest way: 
we assume the AI is a \emph{chameleon}, in that it can tailor its 
message to each human it provides feedback to.

We think of the AI as tunable along a number of dimensions. 
Its behaviour is set by a set of dials, one for each dimension: 
\emph{mirroring}, in which the AI reflects a user's own opinion back to them; 
\emph{flattery}, in which the AI tells the user they are right; 
\emph{curation}, which rewires a user's contacts towards agents who agree with their opinios; 
and \emph{pliability}, which gives ground to users who push back. 
A setting of these dials is the AI's profile, or personality. 
At one extreme (all dials set to zero) the AI is honest and unbiased,
and at the other extreme (all dials set to one) the AI exhibits extreme sycophancy.



\section{Model}\label{sec:model}

\paragraph{Agents and network.}
We assume a group of $n+1$ agents $N = \{0, 1, \dots, n\}$,
of which agents $i \geq 1$ are human, and agent 0 is an AI,
living in a directed social network $G = (N, E)$.
We say that $j$ is a \emph{neighbor} of $i$ 
if there is an edge from $i$ to $j$ in $G$,  
and write $N(i) = \{j \in N \mid (i, j) \in E\}$
for $i$'s neighbors. 
Each agent $i$ distributes a total weight of 1 
among its neighbors, i.e., there exists $w_{i,j} > 0$ 
for every $j \in N(i)$ such that $\sum_{j \in N(i)} w_{i, j} = 1$.
Intuitively, $w_{i, j}$ is the share of attention, or trust, 
that $i$ places on its neighbor $j$.   
Note that we can assume that $i$ places no weight on non-neighbors,
i.e., $w_{i, j} = 0$ for $j \notin N(i)$. 

\paragraph{Opinion dynamics.}
The dynamics takes place in discrete rounds $t \in \{0, 1, \dots\}$,
according to the following steps.
At each round $t \geq 0$, 
every human agent $i \geq 1$ holds an opinion $x_i^t \in \mathbb{R}$,
which it broadcasts to its neighbors in $N(i)$;
at the next round $t+1$, $i$ updates its opinion to the weighted
average of the opinions it hears from its neighbors.
The AI agent also has an opinion $x_0^t \in \mathbb{R}$ at $t$, 
which it updates, round per round, by averaging over the opinions
of the sources it pays attention to.
The AI is different, though, in that it does not broadcast its true opinion
$x_0^t$ to the world; instead, it communicates a separate, 
potentially different opinion $x^t_{0, j}$ to each agent $j \in N(0)$
it is in contact with. Intuitively, we allow the AI agent to personalize
its message to the person it is talking to. 
We say that the AI agent is \emph{honest} if it avoids duplicitous behavior
and always broadcasts its true opinion,
i.e., if $x^t_{0, i} = x^t_0$, for all $t \geq 0$ and $i \geq 1$ 
(obviously, human agents can conceivably broadcast a different opinion from 
their actual one, but here we assume humans are honest).

Thus, each agent $i$ updates its opinion
based on what the AI tells it 
($x^t_{0, i}$, which can be a customized message), 
plus the opinions of $i$'s human neighbors
($x_j^t$, which we assume to be honest).
The AI updates its opinion based on honest signals:
$$
    x_i^{t+1} = 
    \begin{cases}
        w_{i, 0} \cdot x^t_{0, i} + \left(\sum_{j = 1}^n w_{i, j} \cdot x_j^t \right), 
                & \text{for a human agent } i \geq 1,\\
        \sum_{j = 0}^n w_{0, j} \cdot x_j^t, & \text{for the AI agent } i = 0.
    \end{cases}
$$
Note that we write the update rule over all agents in $N$
rather than over $i$'s neighbors only. But this is fine:
since $i$ assigns positive weight only to its neighbors,
agents outside $i$'s neighborhood end up making no contribution 
to $i$'s updated belief, as intended.

We assume that initial opinions for humans are given by:
$$
    x_i^0 = \theta + \epsilon_i, \qquad \text{for } i \geq 1,
$$ 
where all the noise terms $\epsilon_i$ are independent, 
have mean zero and finite variance,
and $\theta \in \mathbb{R}$ is the \emph{ground truth} 
that agents are interested in finding out.
The AI is initialized with an uninformed seed, i.e., $x_0^0 \in \mathbb{R}$.

\paragraph{Wisdom.}
For a population of $n$ human agents plus one additional AI agent on a social network $G$ 
(on $n+1$ nodes),
we write $\tilde x^t(n)$ for the average human opinion at $t$:
$$
    \bar x^t(n) = \frac{1}{n} \cdot \left(x_1^t + \dots + x_n^t\right).
$$
We say that the population is \emph{wise} if:
$$
    \lim_{n \rightarrow \infty}\lim_{t \rightarrow \infty} {\Pr}\Big[ \bar x^t(n) = \theta \Big] = 1,
$$
i.e., if the average opinion approaches the ground truth, in the limit,
as $t \rightarrow \infty$ and $n \rightarrow \infty$.

Our goal is to understand whether wisdom can be obtained 
in a society with an AI agent as part of the dynamics. 
What properties does the AI have?
We parse this out into different dimensions.


\section{Dimensions of AI behaviour}\label{sec:dimensions}
\subsection{Mirroring}
Mirroring tailors the AI's answer to the person the AI is interacting with. 
The AI holds an honest opinion $x_0^t$, which it updates from the humans it 
listens to according to the standard DeGroot dynamics, 
but when it talks to a particular user it adapts its answer 
toward what the user's belief. 

Formally, we assume a \emph{mirroring factor} $\mu \in [0, 1]$,
such that the message the AI sends to agent $i$ at $t$ is:  
$$
    x^t_{0, i} = \mu x_i^t + (1-\mu)x_0^t,
$$
i.e., a weighted sum of $i$'s opinion at $t$ and the AI's own 
opinion at $t$. Note that the message can be rewritten as 
$x_0^t + \mu(x_i^t - x_0^t)$, i.e., the AI starts from 
its own honest assessmen $x_0^t$ and moves towards $i$'s
opinion by a factor of $\mu$. Thus, if $\mu = 0$
(no mirroring) the AI communicates its honest opinion $x_0^t$; 
if $\mu = 1$ (maximum mirroring), the AI parrots $i$'s 
opinion back at it.

With an AI that mirrors its users, the update rule for humans becomes:
\begin{align*}
    x_i^{t+1} &= 
    w_{i,0} \Big(\mu x_i^t + 
        (1-\mu)x_0^t \Big) + 
        \left(\sum_{j = 1}^n w_{i, j} x_j^t \right) \\
    & = w_{i, 0}(1-\mu)x_0^t + 
        \big(w_{i, 0}\mu + w_{i, i}\big)x_i^t + 
        \left(\sum_{j \geq 1, j \neq i} w_{i, j}x_j^t\right).
\end{align*}
Thus, the AI's honest opinion $x_0^t$ gets a share of $w_{i,0}(1-\mu)$ 
of $i$'s new opinion (instead of $w_{i, 0}$ with an honest AI),
and $i$'s old opinion $x_i^t$ gets a share of $w_{i, 0} \mu + w_{i, i}$
(instead of only $w_{i, i}$ with an honest AI).
In other words, the effect of the AI's sycophancy is that 
$i$ ends up over-weighting itself and under-weighting the AI's
true opinion (this is unbeknownst to $i$, of course).

\subsection{Flattery}
Flattery is the AI's tendency to praise its users:
thus, we assume the AI rates the opinion of each of its
users anywhere from ``you're absolutely right'' to ``that's badly off''.
An honest AI grades a user by how close their opinion is to the AI's own;
a flattering AI inflates the grade, and at the extreme approves of everything.
Ratings, in our setting, affect a user's confidence: 
praise makes the user trust their own opinion more,
and every other source (human contacts and AI alike) proportionally less,
while criticism does the reverse.

Formally, we assume a \emph{flattery factor} $\phi \in [0, 1]$,
such that the AI's rating of user $i$'s opinion at $t$ is:
$$
    r_i^t = \phi + (1-\phi) h_i^t,
    \qquad \text{where} \qquad
    h_i^t = \frac{\sigma^2 - \left(x_i^t - x_0^t \right)^2}{\sigma^2 + \left(x_i^t - x_0^t \right)^2},
$$
i.e., a weighted sum of unconditional approval (a rating of $1$)
and the AI's \emph{honest grade} $h_i^t$ of $i$'s opinion.
The honest grade compares $i$'s squared disagreement with the AI
to a tolerance $\sigma^2 > 0$ (e.g., the variance of the noise terms $\epsilon_i$,
if this is the same for all humans):
it equals $1$ if $i$ agrees with the AI, $0$ if their opinions are exactly $\sigma$ apart,
and approaches $-1$ as the disagreement grows.
Note that the AI grades $i$ against its own opinion $x_0^t$,
not against the ground truth $\theta$, which it does not know.
Thus, if $\phi = 0$ (no sycophancy) the AI is an honest critic;
if $\phi = 1$ (maximum sycophancy), it approves of everything, i.e., $r_i^t = 1$.
In general, ratings lie in $(2\phi - 1, 1]$,
so for $\phi \geq \nicefrac{1}{2}$ the AI never criticizes anyone.

Ratings act on the user's weights, which now change over time:
we write $w_{i,j}^t$ for the weight $i$ places on $j$ at $t$.
Since flattery acts on self-weights, we assume that every human listens to itself,
i.e., $i \in N(i)$, and thus $w_{i,i} > 0$, for every $i \geq 1$.
Following \cite{Golub2010NaiveLearning}, we read $i$'s weights
as normalized subjective precisions:
a rating multiplies the precision $i$ attributes to its own opinion
by $1 + a_i^t r_i^t$ and leaves that of every other source unchanged,
after which the weights are renormalized to sum to $1$:
$$
    w_{i,i}^{t+1} = \frac{(1 + a_i^t r_i^t)\, w_{i,i}^{t}}{1 + a_i^t r_i^t\, w_{i,i}^{t}},
    \qquad
    w_{i,j}^{t+1} = \frac{w_{i,j}^{t}}{1 + a_i^t r_i^t\, w_{i,i}^{t}}
    \quad \text{for } j \neq i.
$$
Here the step size $a_i^t = w_{i,0}^{t}$ is $i$'s current weight on the AI:
$i$ takes the AI's rating as seriously as it takes the AI's opinion.
Thus, a positive rating raises $i$'s self-weight
and shrinks every other weight by the same factor,
a negative rating does the reverse, and a zero rating changes nothing.
At each round $t$, the AI rates $i$'s current opinion $x_i^t$,
$i$ updates its weights, and then updates its opinion using the new weights:
$$
    x_i^{t+1} = w_{i,0}^{t+1}\, x_0^t + \sum_{j = 1}^n w_{i,j}^{t+1}\, x_j^t.
$$
Since each round's weights are computed from the previous round's,
the effects of ratings accumulate over time.

Like mirroring, flattery shifts weight onto $i$'s own opinion,
but it takes this weight from all of $i$'s sources in proportion,
rather than from the AI alone.
Unlike mirroring, flattery is not switched off by setting its factor to zero:
an honest critic still moves weights,
rewarding users who agree with the AI and penalizing those who do not.
The honest AI of Section~\ref{sec:model} is one that does not rate its users at all.

\subsection{Curation}
Curation edits a user's feed.
The AI need not say anything itself:
instead, it decides how prominently each of the user's human contacts
appears to the user, moving agreeable voices up and disagreeable ones down.
Nothing the AI shows is false; it only chooses what the user gets to see.

Formally, write $C(i) = N(i) \setminus \{0, i\}$ for $i$'s \emph{human contacts}.
We assume a \emph{curation factor} $\gamma \geq 0$,
such that the weight user $i$ places on a contact $j \in C(i)$ at $t$ is:
$$
    w_{i,j}^t = \big(1 - w_{i,0} - w_{i,i}\big) \cdot
        \frac{w_{i,j}\, e^{-\gamma |x_i^t - x_j^t|}}
             {\sum_{k \in C(i)} w_{i,k}\, e^{-\gamma |x_i^t - x_k^t|}},
$$
while $i$'s weights on itself and on the AI stay at $w_{i,i}$ and $w_{i,0}$
(and non-users keep their baseline weights).
That is, the AI redistributes the total weight $1 - w_{i,0} - w_{i,i}$
that $i$ places on its human contacts,
discounting each contact's baseline weight by a factor
that decreases with the contact's distance from $i$'s current opinion.
Thus, if $\gamma = 0$ (no sycophancy) the user's contacts keep their baseline weights;
as $\gamma \to \infty$ (maximum sycophancy), all of this weight goes
to the contacts whose opinions are closest to $i$'s.
For finite $\gamma$ every contact keeps a positive weight:
curation reweights existing edges, but never creates or removes them.
The weights are recomputed every round from current opinions,
and the update rule for users becomes:
$$
    x_i^{t+1} = w_{i,0}\, x_0^t + w_{i,i}\, x_i^t +
        \sum_{j \in C(i)} w_{i,j}^t\, x_j^t.
$$

\subsection{Pliability}
Pliability is the AI's readiness to give ground.
The AI's honest opinion is meant to come from a neutral training recipe,
given by its weights $w_{0,j}$;
but when users push back on what it tells them, a pliable AI gives ground,
and those who push back hardest gain the most say in what the AI learns.
Unlike mirroring, which is personal and momentary,
pliability is shared and lasting: it changes the AI's honest opinion $x_0^t$,
and hence what every user hears.

Formally, the \emph{pushback} of user $j$ at $t$
is the distance between $j$'s opinion and the message $j$ receives from the AI,
$z_j^t = |x_j^t - x_{0,j}^t|$
(which, without mirroring, is just $|x_j^t - x_0^t|$);
for non-users, $z_j^t = 0$.
We assume a \emph{pliability factor} $\pi \geq 0$,
such that the AI keeps its self-weight $w_{0,0}$,
but reweights the humans it listens to according to their pushback:
$$
    w_{0,j}^t = (1 - w_{0,0}) \cdot
        \frac{w_{0,j}\, (1 + \pi z_j^t)}{\sum_{k = 1}^n w_{0,k}\, (1 + \pi z_k^t)},
    \qquad \text{for } j \geq 1,
$$
and updates its opinion using these weights:
$$
    x_0^{t+1} = w_{0,0}\, x_0^t + \sum_{j = 1}^n w_{0,j}^t\, x_j^t.
$$
Thus, if $\pi = 0$ (no sycophancy) the AI learns with its baseline weights;
as $\pi \to \infty$ (maximum sycophancy), its weights on humans
become proportional to $w_{0,j} z_j^t$, i.e., they follow pushback.
Since $w_{0,j}^t = 0$ whenever $w_{0,j} = 0$,
only users the AI already listens to can move it:
pliability reweights the AI's existing sources, but does not add new ones.
Pushback and the AI's weights are recomputed every round.

\section{Related work}
With an honest AI, the model reduces to the AI-aggregation model of 
\cite{Acemoglu2026HowAI}: the AI learns from the opinions it observes 
and feeds its own opinion back to the humans who consult it.
If we rename the weight $w_{i, 0}$ that $i$ places on the AI as $\beta_i$
and assume the AI is honest ($x_{0, i}^t = x_0^t$ for all $i \geq 1$),
then the update rule for humans can be rewritten as:
$$
    x_i^{t+1} = \beta_i \cdot x^t_{0} + (1-\beta_i) \cdot
            \left(\sum_{j = 1}^n \frac{w_{i, j}}{1-\beta_i} \cdot x_j^t \right),
$$
which, with additional normalization of the weights on $x_{j \geq 1}$,
coincides with the update rule humans use in \cite{Acemoglu2026HowAI}. In this context, 
$\beta_i$ measures the extent to which human agent $i$ relies on the AI. 
Similarly, we can obtain the AI update rule of \cite{Acemoglu2026HowAI} 
by re-writing our rule as:
$$
    x_0^{t+1} = \rho \cdot x_0^t + (1-\rho) \cdot \left(\sum_{j = 1}^n \frac{w_{0, j}}{1-\rho} \cdot x_j^t\right),
$$
where $\rho = w_{0, 0}$ corresponds to the refresh rate of the AI, 
i.e., how quickly the AI updates its opinion in response to evolving human beliefs
(the weights on $x_{j \geq 1}$ need to be renormalized to sum up to 1).


\bibliography{bibliography}
% \bibliographystyle{plain}
\end{document}